\documentclass[11pt,a4paper]{article}
\usepackage[T1]{fontenc}
\usepackage[utf8]{inputenc}
\usepackage[brazilian]{babel}
\usepackage{lmodern,microtype}
\usepackage[a4paper,margin=2.4cm]{geometry}
\usepackage{enumitem,booktabs,longtable,array,amssymb}
\usepackage{xcolor,hyperref,xurl}
\hypersetup{colorlinks=true,urlcolor=blue!55!black,linkcolor=blue!55!black}
\setlength{\parindent}{0pt}
\setlength{\parskip}{5pt}
\setlist[enumerate]{leftmargin=2em,itemsep=4pt}
\begin{document}
\begin{center}
{\LARGE\bfseries Temas das provas didáticas de Matemática do IFMT\par}
\vspace{4pt}
{\large Editais 166/2023, 029/2024 e 003/2026\par}
\vspace{6pt}
{\small Levantamento documental, revisão de 4 de novembro de 2026\par}
\end{center}
\tableofcontents
\newpage

\section{Escopo e método de conferência}
Este arquivo reúne os temas de Matemática previstos para o sorteio da prova de desempenho didático nos dois editais anteriores cujo anexo específico pôde ser consultado integralmente. Para 2026, identifica o anexo obrigatório e a situação da verificação. \textbf{Não se atribui a 2026 uma lista inferida de programas anteriores.}

\textbf{Distinção essencial:} o \emph{Anexo III}, em cada edital, trata de conteúdos para prova \textbf{objetiva}; os temas didáticos têm anexo próprio: \textbf{Anexo X} em 2023/2024 e \textbf{Anexo XVI} em 2026. Os cinco temas são possibilidades de sorteio, não cinco aulas realizadas pelo mesmo candidato.

\section{Edital 166/2023 --- Anexo X: Matemática}
Documento-base: anexos consolidados do Edital 166/2023, com retificações. Fonte oficial:
\url{https://seletivo.ifmt.edu.br/edital/arquivos/889/}.

\begin{enumerate}
\item Aplicações de derivadas: máximos e mínimos locais.
\item Integral definida e suas aplicações.
\item Limites de funções reais de uma variável real.
\item Teorema do Valor Médio e Teorema de Taylor.
\item Transformações lineares e matrizes.
\end{enumerate}

\subsection{Possíveis recortes de preparação (interpretação pedagógica)}
\begin{description}[leftmargin=3.2cm,style=nextline]
\item[Derivadas] Pontos críticos, teste de sinais, diferenciação entre extremos locais e globais, problemas de otimização.
\item[Integral] Soma de Riemann, Teorema Fundamental do Cálculo, interpretação de acumulação e unidades.
\item[Limites] Limites laterais, indeterminação removível, demonstração pontual por definição.
\item[Valor Médio/Taylor] Hipóteses de Rolle e TVM, aproximação polinomial e estimativas de erro.
\item[Álgebra Linear] Bases, representação matricial, transformações e composição.
\end{description}
Os subitens deste parágrafo \emph{não são uma ampliação literal dos pontos do edital}; são ideias autorais para preparar aulas e exercícios.

\section{Edital 029/2024 --- Anexo X \textbf{retificado}: Matemática}
A referência correta é o \textbf{Edital Retificador 01/2024}, de 30/08/2024, que retificou e complementou o Anexo X. O texto anterior de abertura não deve prevalecer sobre a retificação.

Fonte oficial:
\url{https://seletivo.ifmt.edu.br/edital/arquivos/1837/}.
\begin{enumerate}
\item \textbf{Geometria plana e espacial:} noções fundamentais; congruências de figuras planas; propriedades do triângulo; Teorema de Thales; áreas e perímetros de figuras planas; figuras no espaço: posições relativas de retas e planos; volume e área superficial de sólidos; unidades de medida do Sistema Internacional.
\item \textbf{Limites e continuidade} de funções reais de uma variável real.
\item \textbf{Técnicas de integração}; integral definida e aplicações.
\item \textbf{Equações diferenciais:} equações diferenciais de primeira e segunda ordem com coeficientes constantes.
\item \textbf{Espaços vetoriais e transformações lineares:} espaços vetoriais de dimensão finita; transformações lineares; Teorema do Núcleo e da Imagem.
\end{enumerate}
\textbf{Observação:} o primeiro ponto é muito amplo. Uma aula de duração limitada exige \emph{recorte explícito} e coerência com eventual orientação do edital e sorteio; nunca apresentar uma pequena parte como cobertura de todo o assunto.

\section{Edital 003/2026 --- situação dos temas didáticos}
A publicação de abertura do Edital 003/2026 estabelece que os temas sorteáveis para cada área constam do \textbf{Anexo XVI}. O edital de 2026 prevê uma aula de \textbf{30 a 35 minutos} em sua prova didática, respeitadas retificações.

\textbf{Situação documental em 04/11/2026:} não foi possível obter e conferir integralmente, por fonte acessível e identificável, o \textbf{Anexo XVI da área de Matemática}. Por isso, os temas de 2026 não são reproduzidos e não devem ser deduzidos dos editais 166/2023 ou 029/2024.

Página oficial de acompanhamento:
\url{https://seletivo.ifmt.edu.br/edital/visualizar/173/}.

Assim que o arquivo oficial for acessível, inserir aqui a transcrição fiel dos pontos, com identificação de versão, data e eventual retificação; conferir também a existência de sorteio de tema já publicado.

\section{Comparação de incidência temática}
\begin{longtable}{p{.39\textwidth}p{.23\textwidth}p{.23\textwidth}}
\toprule
\textbf{Eixo}&\textbf{166/2023}&\textbf{029/2024 (ret.)}\\\midrule
\endhead
Análise e Cálculo diferencial & Derivadas, limites, TVM, Taylor & Limites e continuidade\\
Integrais & Integral definida e aplicações & Técnicas de integração e aplicações\\
Álgebra Linear & Transformações e matrizes & Espaços vetoriais, transformações, núcleo e imagem\\
Geometria & Não consta como tema autônomo & Geometria plana e espacial\\
EDOs & Não consta como tema autônomo & 1ª e 2ª ordem com coeficientes constantes\\
\bottomrule
\end{longtable}
\textbf{Leitura estratégica:} dominar integrais, limites e Álgebra Linear é útil para ambos os programas; Geometria e EDOs surgem explicitamente em 2024. A tabela registra apenas a \emph{presença de pontos didáticos}, sem prever probabilidade de cobrança futura.

\section{Checklist de conferência antes da prova}
\begin{enumerate}[label={$\square$}]
\item Verifiquei o número e a versão do anexo de \textbf{temas didáticos}, não apenas o programa objetivo.
\item Identifiquei precisamente os pontos da \textbf{área Matemática}, e não de outra área.
\item Li as retificações antes de selecionar o tema.
\item Defini turma, pré-requisitos, objetivo, método, avaliação e bibliografia compatíveis.
\item Preparei resoluções previstas e exemplos com passos verificáveis.
\item Treinei o tempo do edital vigente com conclusão e avaliação, sem encerramento abrupto.
\end{enumerate}
\section*{Referências}
\begin{enumerate}
\item IFMT, Edital 166/2023, anexos consolidados (Anexo X). \url{https://seletivo.ifmt.edu.br/edital/arquivos/889/}.
\item IFMT, Edital 029/2024, Retificador 01/2024 (Anexo X retificado). \url{https://seletivo.ifmt.edu.br/edital/arquivos/1837/}.
\item IFMT, Edital 003/2026, página de acompanhamento, Anexo XVI indicado no edital. \url{https://seletivo.ifmt.edu.br/edital/visualizar/173/}.
\end{enumerate}
\end{document}
